% Propietario conservado: /Users/ruben/Documents/New project/output/EXTREMA_MEDIA_RAZON_RECIPROCIDAD_ES_EN_20260918/ARTICULO/ES/source/libro/03_refinamiento.tex
\chapter{Conservative refinement of unity and carry memory}
\label{x_reciprocidad:lib03:refinamiento}

Dimensional completion and the record of operations admit an explicit
arithmetic composition. The APP--TRIT--TPK architecture first provides
the enriched state, the oriented readings, and the
completions of unity. From these outputs we construct an
affine transport that preserves the normalized sum and permits recovery, by
Euclidean division, of every finite word of canonical carries. Its limit
has cylinders of controlled diameter; boundary marks distinguish
the histories that a scalar evaluation identifies.

The central block, its concatenations, and the partition
\(1000=729+243+27+1\) are antecedents of this construction. The composition
of these objects with the affine transport, the residue inverse, and its
prolongation constitutes the formalization developed in this section.
The map acts on a declared arithmetic domain. Each realization
on a particular enriched history also retains the
prospective rule that produces its carries.

\section{Relation between quadratic norm and cubic completion}
\label{x_reciprocidad:lib03:completacion}

Let \(b\geq1\) be an integer and \(B=b+1\geq2\). Define
\begin{equation}
 p_b=b^2-b+1,\qquad q_b=B^2-p_b=3b.
 \label{x_reciprocidad:lib03:pb}
\end{equation}

\begin{proposition}[Transfer of the completion]
\label{x_reciprocidad:lib03:prop-completacion}
The identities
\begin{equation}
 Bp_b=b^3+1,\qquad Bq_b+1=B^3-b^3
 \label{x_reciprocidad:lib03:cuadratica-cubica}
\end{equation}
determine the transformation
\begin{equation}
 (p_b,q_b)\longmapsto(Bp_b-1,Bq_b+1)
       =(b^3,B^3-b^3).
 \label{x_reciprocidad:lib03:transferencia-cubica}
\end{equation}
The sum changes from \(B^2\) to \(B^3\), while the sum of the
coordinates divided by their respective scales remains equal to one.
\end{proposition}

\begin{proof}
The factorization \(b^3+1=(b+1)(b^2-b+1)\) yields the first
identity. The second follows from
\((b+1)^3-b^3=3b^2+3b+1\). Adding them proves the change of scale and,
after division by \(B^3\), the conservation of normalized unity.
\end{proof}

Para \(b=9\),
\begin{equation}
 (73,27)\longmapsto(10\cdot73-1,10\cdot27+1)=(729,271).
 \label{x_reciprocidad:lib03:caso-729}
\end{equation}
The complement preserves the decomposition \(271=243+27+1\). In the
preceding completion, the last contribution corresponds to the added
apex. The correction \((-1,+1)\) has zero sum, and multiplication
by ten transforms the scale of one hundred into that of one thousand. For the complete
family, the transport of proportions is
\begin{equation}
 \left(\frac{p_b}{B^2},\frac{q_b}{B^2}\right)
 \longmapsto
 \left(\frac{p_b}{B^2}-\frac1{B^3},
       \frac{q_b}{B^2}+\frac1{B^3}\right).
 \label{x_reciprocidad:lib03:fracciones-completacion}
\end{equation}
The quantity preserved by this identity is a linear sum. Quadratic
forms and physical quantities have their own readers, which will
subsequently be composed with this transport.

\section{Concatenations of the central block and restoration of unity}
\label{x_reciprocidad:lib03:centro}

For \(b\geq2\), the digits \(b-2\) and \(2\) belong to the alphabet of
base \(B=b+1\). Consider the algebraic family
\begin{equation}
 B_c^{(b)}=\begin{pmatrix}b-2&2\\2&b-2\end{pmatrix},\qquad
 C_B(u,v)=Bu+v.
 \label{x_reciprocidad:lib03:bloque-central}
\end{equation}
The case \(b=9\) recovers the documented central block
\(\bigl(\begin{smallmatrix}7&2\\2&7\end{smallmatrix}\bigr)\), with
eigenvalues \(9\) and \(5\). Its concatenations satisfy
\(72+27=77+22=99\). Concatenation, spectrum, and central position are
distinct readings of the same antecedent and retain their domains.

Applying the reader to the rows of \eqref{x_reciprocidad:lib03:bloque-central} gives
\begin{align}
 C_B(b-2,2)&=B(b-2)+2=b(b-1)=p_b-1,\\
 C_B(2,b-2)&=2B+(b-2)=3b=q_b.
 \label{x_reciprocidad:lib03:concatenaciones}
\end{align}
The sum is \(B^2-1\). Restoring one unit in the first
coordinate and performing the subsequent transfer yields
\begin{equation}
 (p_b-1,3b)\xrightarrow{\ +(1,0)\ }(p_b,3b)
 \xrightarrow{\ L_1\ }(b^3,B^3-b^3).
 \label{x_reciprocidad:lib03:cadena-restauracion}
\end{equation}
In particular,
\begin{equation}
 (72,27)\longmapsto(73,27)\longmapsto(729,271).
 \label{x_reciprocidad:lib03:cadena-central}
\end{equation}
The first arrow increases the sum by one. The second changes the scale
and redistributes one unit through a zero-sum correction. To represent
the contribution explicitly before incorporating it, the triple
\((p_b-1,3b,1)\) already has sum \(B^2\), and the map
\((x,y,r)\mapsto(x+r,y,0)\) preserves that sum. Its inversion requires
retaining the datum \(r=1\): on arbitrary triples the map is not
injective.

Restitution is thus recorded without attributing to it a closed
conservation of the pair before incorporating the unit. The family in \(b\)
algebraically extends the received central case; the selection of
atlas states remains that established by APP--TRIT--TPK.
The composition \eqref{x_reciprocidad:lib03:cadena-restauracion} determines a carry
equal to one in that transition, leaving open the prospective
specification of subsequent operations.

\section{Affine transport and the positivity domain}
\label{x_reciprocidad:lib03:afin}

\begin{definition}[Transport with carry]
\label{x_reciprocidad:lib03:def-transporte}
For an integer base \(B\geq2\) and \(c\in\mathbb Z\), define
\begin{equation}
 L_c(x,y)=(Bx-c,By+c).
 \label{x_reciprocidad:lib03:lc}
\end{equation}
The integer \(c\) specifies the oriented transfer between the two
coordinates. A positive sign reduces the first relative to \(Bx\)
and increases the second by the same amount.
\end{definition}

\begin{proposition}[Inverse and positivity]
\label{x_reciprocidad:lib03:prop-inversa}
In the subsequent affine realization \(\mathbb R^2\), with \(B,c\) fixed,
\begin{equation}
 L_c^{-1}(X,Y)=\left(\frac{X+c}{B},\frac{Y-c}{B}\right),\qquad
 X+Y=B(x+y).
 \label{x_reciprocidad:lib03:inversa-afin}
\end{equation}
The preimage is integral exactly when \(X+c\) and \(Y-c\) are
divisible by \(B\). For \(x,y\geq0\), the image is nonnegative
exactly when
\begin{equation}
 -By\leq c\leq Bx.
 \label{x_reciprocidad:lib03:positividad}
\end{equation}
\end{proposition}

\begin{proof}
Direct substitution proves both identities in
\eqref{x_reciprocidad:lib03:inversa-afin} and the integrality criterion. The
inequalities \(Bx-c\geq0\), \(By+c\geq0\) are equivalent to
\eqref{x_reciprocidad:lib03:positividad}.
\end{proof}

For the canonical alphabet \(c\in\{0,\ldots,B-1\}\), all
prolongations are admissible if \(x\geq1\), \(y\geq0\). At the
boundary \(x=0\), nonnegativity permits only \(c=0\).
The positive domain and the general integer domain are thus
distinguished by an explicit condition.

\Needspace{12\baselineskip}
\section{Canonical recovery through the residue}
\label{x_reciprocidad:lib03:residuo}

\begin{theorem}[Canonical inverse by residue]
\label{x_reciprocidad:lib03:teo-residuo}
Let \(M\geq1\) be an integer and let \(X,Y\geq0\) be integers such that
\(X+Y=BM\). There is a unique admissible triple \((x,y,c)\) satisfying
\[
 x+y=M,\qquad 0\leq c\leq B-1,\qquad L_c(x,y)=(X,Y).
\]
It is obtained from
\begin{equation}
 c=Y\bmod B=(-X)\bmod B,\qquad
 y=\frac{Y-c}{B},\qquad x=\frac{X+c}{B}.
 \label{x_reciprocidad:lib03:residuo-inversor}
\end{equation}
\end{theorem}

\begin{proof}
Every preimage satisfies \(Y\equiv c\pmod B\). The representative in
\(\{0,\ldots,B-1\}\) is unique. The congruence \(X+Y\equiv0\pmod B\)
implies the divisibility of \(X+c\). Euclidean division of \(Y\)
gives \(y\geq0\), and \(X+c\geq0\) gives \(x\geq0\). If
\(c>0\), the case \(x=0\) would require \(X=-c<0\), so the
resulting triple belongs to the admissible domain. Its sum is \(M\), and
substitution into \eqref{x_reciprocidad:lib03:lc} recovers \((X,Y)\). Uniqueness
of the residue and the two quotients proves uniqueness of the triple.
\end{proof}

The integer pair and its scale contain the last canonical carry. Its
recovery follows from the residue and does not require retaining another
independent copy of that carry. There are \(M+1\) triples with \(c=0\) and
\(M\) triples in each of the remaining \(B-1\) classes: the total
is \(BM+1\), equal to the number of nonnegative pairs with sum \(BM\).

For a general integer carry, write
\(c=r+B\ell\), with \(0\leq r\leq B-1\). The residue determines
\(r\), whereas complete reconstruction also uses the
lift \(\ell\):
\begin{equation}
 y=\frac{Y-r}{B}-\ell,\qquad
 x=\frac{X+r}{B}+\ell.
 \label{x_reciprocidad:lib03:elevacion}
\end{equation}
Retaining the lift distinguishes the signed carry from its
residual reduction.

\section{Variable iteration and complete finite memory}
\label{x_reciprocidad:lib03:memoria}

Fix \((x_0,y_0)\in\mathbb Z^2\), with sum \(M_0>0\), and a
sequence of operations \(c_1,c_2,\ldots\). Define
\begin{equation}
 (x_n,y_n)=L_{c_n}\cdots L_{c_1}(x_0,y_0),\qquad
 C_n=\sum_{j=1}^n c_jB^{n-j},\qquad C_0=0.
 \label{x_reciprocidad:lib03:registro-acumulado}
\end{equation}
The carries are supplied by the corresponding operational rule; the
limit value does not participate in their selection.

\begin{theorem}[Exact accumulated record]
\label{x_reciprocidad:lib03:teo-acumulacion}
For every \(n\geq0\),
\begin{equation}
 \begin{aligned}
 x_n&=B^nx_0-C_n,& y_n&=B^ny_0+C_n,\\
 x_n+y_n&=B^nM_0,& C_{n+1}&=BC_n+c_{n+1}.
 \end{aligned}
 \label{x_reciprocidad:lib03:iteracion-exacta}
\end{equation}
\end{theorem}

\begin{proof}
For \(n=0\), the identities are immediate. If they hold at
\(n\), applying \(L_{c_{n+1}}\) multiplies all
previous contributions by \(B\) and adds \(c_{n+1}\) to the accumulated carry.
This gives the recurrence and, by substitution, both coordinates of the
next stage. Their sum eliminates \(C_n\).
\end{proof}

When the carries are canonical,
\begin{equation}
 0\leq C_n\leq B^n-1.
 \label{x_reciprocidad:lib03:intervalo-acumulado}
\end{equation}
In particular, \(x_0\geq1\), \(y_0\geq0\) guarantee
\(x_n\geq B^n(x_0-1)+1\) and \(y_n\geq0\) at every stage.

\begin{corollary}[Recovery of the complete word]
\label{x_reciprocidad:lib03:cor-palabra}
Given \(B,n\) and the final pair, a finite canonical history
recovers its input and all its carries through
\begin{align}
 C_n&=y_n\bmod B^n,&
 y_0&=\left\lfloor\frac{y_n}{B^n}\right\rfloor,&
 x_0&=\frac{x_n+C_n}{B^n},
 \label{x_reciprocidad:lib03:recuperacion-inicial}\\
 c_j&=\left\lfloor\frac{C_n}{B^{n-j}}\right\rfloor\bmod B,
 &&1\leq j\leq n.
 \label{x_reciprocidad:lib03:recuperacion-digitos}
\end{align}
\end{corollary}

\begin{proof}
The bound \eqref{x_reciprocidad:lib03:intervalo-acumulado} makes
\(y_n=B^ny_0+C_n\) a Euclidean division. Its quotient and residue
give \(y_0,C_n\); the first coordinate gives \(x_0\).
Successive divisions of \(C_n\) by powers of \(B\) recover
the digits. Equivalently, apply
\eqref{x_reciprocidad:lib03:residuo-inversor} from the last step to the first.
\end{proof}

The depth retains the leading zeros of the word, even though they
do not change the integer \(C_n\). If \(M_0\) is stored instead of
\(n\), the exact equality \((x_n+y_n)/M_0=B^n\) recovers the
depth. Normalizing the pair to sum to one removes that scale counter
unless it is retained by another reading.

\section{Block composition and compatibility with the partition}
\label{x_reciprocidad:lib03:composicion}

A block of \(n\) operations, with accumulated carry \(C\), acts by
\begin{equation}
 L_C^{[n]}(x,y)=(B^nx-C,B^ny+C).
 \label{x_reciprocidad:lib03:bloque}
\end{equation}
The composition of two blocks satisfies
\begin{equation}
 L_D^{[m]}\circ L_C^{[n]}=L_{B^mC+D}^{[n+m]}.
 \label{x_reciprocidad:lib03:ley-bloques}
\end{equation}
Indeed, substituting the first image into the second multiplies its
accumulated carry by \(B^m\) and adds \(D\). The law
\[
 (n,C),(m,D)\longmapsto(n+m,B^mC+D)
\]
is associative by composition of maps. It expresses the contribution
of the first block at the scale of the second and preserves the output when
the same ordered events are regrouped. Adding new events changes
the accumulated carry according to this same law.

\section{Convergence and quantitative control of the reading}
\label{x_reciprocidad:lib03:limite}

Write \(M_n=B^nM_0\) and
\begin{equation}
 f_n=\frac{x_n}{M_n},\quad g_n=\frac{y_n}{M_n},\quad
 \theta_n=\frac{C_n}{B^n}=\sum_{j=1}^n\frac{c_j}{B^j}.
 \label{x_reciprocidad:lib03:lecturas-normalizadas}
\end{equation}
The exact identities are
\begin{equation}
 f_n=\frac{x_0}{M_0}-\frac{\theta_n}{M_0},\qquad
 g_n=\frac{y_0}{M_0}+\frac{\theta_n}{M_0},\qquad f_n+g_n=1.
 \label{x_reciprocidad:lib03:unidad-normalizada}
\end{equation}

\begin{theorem}[Convergence of bounded carries]
\label{x_reciprocidad:lib03:teo-limite}
If \(|c_j|\leq C\) for every \(j\), the series
\(\theta=\sum_{j\geq1}c_jB^{-j}\) converges absolutely. The
readings \(f_n,g_n\) converge to \(f,g\), with \(f+g=1\), and
\begin{equation}
 |f-f_n|=|g-g_n|\leq\frac{C}{M_0(B-1)B^n}.
 \label{x_reciprocidad:lib03:cota-cola}
\end{equation}
\end{theorem}

\begin{proof}
The absolute tail satisfies
\[
 \sum_{j>n}\frac{|c_j|}{B^j}
 \leq C\sum_{j>n}B^{-j}=\frac{C}{(B-1)B^n}.
\]
Convergence of the geometric series and
\eqref{x_reciprocidad:lib03:unidad-normalizada} prove all the assertions.
\end{proof}

For carries of both signs, a sufficient condition for positivity
at every stage is \(x_0,y_0\geq C/(B-1)\). For the canonical
alphabet, \(x_0\geq1,y_0\geq0\) suffices. In this second case,
\(0\leq\theta\leq1\), \(f_n\) decreases and \(g_n\) increases, and
\begin{equation}
 0\leq f_n-f=g-g_n\leq\frac1{M_0B^n}.
 \label{x_reciprocidad:lib03:cota-canonica}
\end{equation}
The sum remains fixed while the individual proportions evolve.
The finite history remains recoverable when the integer
coordinates and the scale are retained.

\section{Compatible cylinders and boundary multiplicity}
\label{x_reciprocidad:lib03:cilindros}

Consider the full canonical alphabet and an input
\(x_0\geq1,y_0\geq0\). A prefix of depth \(n\), with
accumulated carry \(C_n\), admits exactly the limit values
\begin{equation}
 I_n(C_n)=\left[\frac{C_n}{B^n},\frac{C_n+1}{B^n}\right]
 \label{x_reciprocidad:lib03:cilindro}
\end{equation}
of the reading \(\theta\).

\begin{theorem}[Nested cylinders and distinction of histories]
\label{x_reciprocidad:lib03:teo-cilindros}
The \(B\) child cylinders of \(I_n(C_n)\) cover the parent and
overlap only at endpoints. The cylinders of an infinite history
have diameter \(B^{-n}\) and a single-point intersection.
Two different canonical histories produce the same value if and only
if, at their first differing position, their digits differ by one and the
subsequent tails are, respectively, all \(B-1\) for the smaller
digit and all zero for the larger.
\end{theorem}

\begin{proof}
The child associated with \(c\) has endpoints
\((BC_n+c)/B^{n+1}\) and \((BC_n+c+1)/B^{n+1}\).
These expressions prove inclusion, coverage and overlap only
at endpoints. The tail of an expansion ranges over the interval
\([0,B^{-n}]\); successively choosing one of the children produces an
expansion of each point in that interval. The nested closed
intervals have diameters tending to zero, which proves existence
and uniqueness of the intersection for each history.

Let \(k\) be the first differing position and let \(d_k>c_k\). If the
readings coincide,
\[
 \frac{d_k-c_k}{B^k}
 =\sum_{j>k}\frac{c_j-d_j}{B^j}
 \leq\sum_{j>k}\frac{B-1}{B^j}=\frac1{B^k}.
\]
The positive integer difference \(d_k-c_k\) must equal one. Equality
in the tail bound requires \(c_j=B-1\), \(d_j=0\) for every
\(j>k\). Conversely, those tails differ by exactly
\(B^{-k}\), which compensates for the first digit. This proves the
complete characterization.
\end{proof}

The image of the cylinder for \(f\) has diameter
\(1/(M_0B^n)\), matching \eqref{x_reciprocidad:lib03:cota-canonica}.
The histories \((1,0,0,\ldots)\) and
\((0,B-1,B-1,\ldots)\) both produce \(\theta=1/B\), while retaining
distinct prefixes and finite integer states. A boundary mark
retains both provenances. A convention excluding tails eventually
constant at \(B-1\) selects one expansion for \(\theta<1\);
the endpoint \(\theta=1\) is retained separately, since in the
fractional alphabet it corresponds to the constant tail \(B-1\).

The history of operations, the cylinder system, and the limit value are
objects related by these maps. The uniqueness of a
history's reading coexists with multiplicity of histories at
boundary values.

\section{Signed carries and prospective specification}
\label{x_reciprocidad:lib03:firmados}

Allowing carries outside the canonical alphabet requires retaining their
lifts. Already in two steps, the words
\[
 (c_1,c_2)=(1,-1),\qquad(d_1,d_2)=(0,B-1)
\]
produce \(C_2=B-1\) and the same final pair. Both are bounded by
\(B-1\) and can preserve positivity, for example starting from
\((73,27)\) in base ten. The signed word, or the lifts in
\eqref{x_reciprocidad:lib03:elevacion}, distinguishes that information. The convergence
theorem remains valid for both words; the canonical inverse,
however, has the domain specified in its statement.

The transition \eqref{x_reciprocidad:lib03:caso-729} fixes \(c_1=1\). An additional
rule imposing \(c_j=1\) for every \(j\) gives
\begin{equation}
 \theta_n=\frac{1-B^{-n}}{B-1},\qquad
 \theta=\frac1{B-1}.
 \label{x_reciprocidad:lib03:acarreo-constante}
\end{equation}
In base ten, starting from \((73,27)\), the sequence is
\[
 (73,27)\longmapsto(729,271)\longmapsto(7289,2711)
 \longmapsto(72889,27111),
\]
and the first fraction converges to
\begin{equation}
 \frac{73}{100}-\frac1{900}=\frac{164}{225}
 =0{,}728888\ldots.
 \label{x_reciprocidad:lib03:limite-constante}
\end{equation}
The conclusion belongs to that additional constant-carry rule.
With only \(c_1=1\) and nonnegative canonical prolongations, the
first-step cylinder yields exactly
\begin{equation}
 \frac{b^3-1}{B^3}\leq f\leq\frac{b^3}{B^3}.
 \label{x_reciprocidad:lib03:primer-cilindro}
\end{equation}
In base ten it is \([0{,}728,0{,}729]\). Changing the sign or the alphabet
changes this family and requires preserving the corresponding
additional memory. The presence of \(729\) in the completion determines
the written arithmetic transition; identification with another reader of
constants also uses its effective generation rule and its
intertwining.

\section{Operator realization and transport amplitudes}
\label{x_reciprocidad:lib03:operativo}

The resulting procedure comprises updating through
\eqref{x_reciprocidad:lib03:lc}, local residue checking, reconstruction of
every canonical word through \eqref{x_reciprocidad:lib03:recuperacion-digitos},
control of the limit reading through \eqref{x_reciprocidad:lib03:cota-cola}, and
regrouping of events through \eqref{x_reciprocidad:lib03:ley-bloques}. Each
operation retains the order and scale counter used by its
inverse.

Conservation of the sum and quadratic conservation have different
forms. For a normalized transfer
\(\delta=c/(BM)\),
\begin{equation}
 (f-\delta)^2+(g+\delta)^2-(f^2+g^2)
 =2\delta(g-f)+2\delta^2.
 \label{x_reciprocidad:lib03:variacion-cuadratica}
\end{equation}
This expression determines the quadratic variation and is generally
nonzero. The connection with an isometry will be obtained on the
amplitudes of the admissible labels, through the residue bijection of
\cref{x_reciprocidad:lib05:levantamiento}. Arithmetic transport and its
unitary realization are thus composed while retaining their respective
quantities, rather than being identified through a scalar coincidence.
